{{دیوان
| عنوان = بلالؓ (لِکھّا ہے ایک مغربیِ حق شناس نے)
| ماخذ = ویکی ماخذ
| ماخذ_لنک = https://ur.wikisource.org/wiki/%D8%A8%D9%84%D8%A7%D9%84_%28%D9%84%D9%90%DA%A9%DA%BE%D9%91%D8%A7_%DB%81%DB%92_%D8%A7%DB%8C%DA%A9_%D9%85%D8%BA%D8%B1%D8%A8%DB%8C%D9%90_%D8%AD%D9%82_%D8%B4%D9%86%D8%A7%D8%B3_%D9%86%DB%92%29
| مقابلہ = اقبال سائبر لائبریری (اقبال اکادمی پاکستان)
| مقابلہ_لنک = https://www.iqbalcyberlibrary.net
}}
<poem>
لِکھّا ہے ایک مغربیِ حق شناس نے
اہلِ قلم میں جس کا بہت احترام تھا

جولاں گہِ سکندرِ رومی تھا ایشیا
گردُوں سے بھی بلند تر اُس کا مقام تھا

تاریخ کہہ رہی ہے کہ رومی کے سامنے
دعویٰ کِیا جو پورس و دارا نے، خام تھا

دنیا کے اُس شہنشہِ انجم سپاہ کو
حیرت سے دیکھتا فلکِ نیل فام تھا

آج ایشیا میں اُس کو کوئی جانتا نہیں
تاریخ دان بھی اُسے پہچانتا نہیں

لیکن بِلالؓ، وہ حبَشی زادۂ حقیر
فطرت تھی جس کی نُورِ نبوّت سے مُستنیر

جس کا امیں ازل سے ہُوا سینۂ بِلالؓ
محکوم اُس صدا کے ہیں شاہنشہ و فقیر

ہوتا ہے جس سے اسوَد و احمر میں اختلاط
کرتی ہے جو غریب کو ہم پہلوئے امیر

ہے تازہ آج تک وہ نوائے جگر گداز
صدیوں سے سُن رہا ہے جسے گوشِ چرخِ پیر

اقبالؔ! کس کے عشق کا یہ فیضِ عام ہے
رومی فنا ہُوا، حبَشی کو دوام ہے
</poem>
